% GENERATED FILE. Edit canonical lessons, then run npm run generate:books.
% canonical-source-hash: fnv1a64:e343aac24d597b8b
% canonical-lessons: HI-C265-patang, HI-C265-jhula, HI-C265-parda, HI-C265-dari, HI-C265-qalin

\chapter{A kite, A swing, A curtain, A cotton rug, A carpet}
\label{ch:hi-a-kite-a-swing-a-curtain-a-cotton-rug-a-carpet}
\hlchaptermodality{265}

\begin{chapteropening}
\textbf{By the end of this chapter:} \emph{I can say a kite, a swing, a curtain, a cotton rug and a carpet.}

Complete the last lesson of chapter 265: I can say a kite, a swing, a curtain, a cotton rug and a carpet.
\end{chapteropening}

\section[\texorpdfstring{\hi{पतंग} (patang)}{पतंग (patang)}]{\hi{पतंग} (patang) — a kite}

\label{lesson:HI-C265-patang}



\begin{quote}
Before the new one: say the Hindi for a toy, then the Hindi for a doll.
\end{quote}

\subsection*{You'll want to know: \hi{पतंग}}
\textbf{\hi{पतंग}} — \emph{patang} — \textquotedblleft{}a kite\textquotedblright{}.

\begin{cousinweb}
\textbf{\hi{पतंग}} (\emph{patang}) is usually traced to Sanskrit \emph{pataṅga}, \textquotedblleft{}bird; winged insect\textquotedblright{}. CDIAL lists that word, with Hindi \emph{patiṅgā}, \textquotedblleft{}grasshopper\textquotedblright{}, as its inherited form (CDIAL 7721), but not \emph{patang}, \textquotedblleft{}kite\textquotedblright{}.
\end{cousinweb}

\begin{grammarlens}[title={What you've built}]
One more word for this chapter.
\end{grammarlens}

\subsection*{Your turn}
\begin{itemize}
\raggedright
  \item \emph{Say it:} \emph{patang}
  \item \emph{Say it:} \emph{patang}, once more
  \item \emph{Say it:} say \emph{patang}
  \item \emph{Recall:} say the Hindi for a toy, then the Hindi for a doll, then say \emph{patang} again
\end{itemize}

\begin{tcolorbox}[breakable,colback=teal!4,colframe=teal!35!black,title={Before you move on}]
What does \hi{पतंग} mean? (\textquotedblleft{}A kite\textquotedblright{}.) Say it once more.
\end{tcolorbox}
\section[\texorpdfstring{\hi{झूला} (jhūlā)}{झूला (jhūlā)}]{\hi{झूला} (jhūlā) — a swing}

\label{lesson:HI-C265-jhula}



\begin{quote}
Before the new one: say the Hindi for a doll, then the Hindi for a kite.
\end{quote}

\subsection*{You'll want to know: \hi{झूला}}
\textbf{\hi{झूला}} — \emph{jhūlā} — \textquotedblleft{}a swing\textquotedblright{}.

\begin{cousinweb}
\textbf{\hi{झूला}} (\emph{jhūlā}) is a tadbhava, a word inherited and reshaped by sound change. It comes from \emph{*jhulyati}, \textquotedblleft{}swings\textquotedblright{}, a form CDIAL reconstructs from its descendants (CDIAL 5406). \emph{jhūlā} is the noun of \emph{jhūlnā}, \textquotedblleft{}to swing\textquotedblright{}.
\end{cousinweb}

\begin{grammarlens}[title={What you've built}]
One more word for this chapter.
\end{grammarlens}

\subsection*{Your turn}
\begin{itemize}
\raggedright
  \item \emph{Say it:} \emph{jhūlā}
  \item \emph{Say it:} \emph{jhūlā}, once more
  \item \emph{Say it:} say \emph{jhūlā}
  \item \emph{Recall:} say the Hindi for a doll, then the Hindi for a kite, then say \emph{jhūlā} again
\end{itemize}

\begin{tcolorbox}[breakable,colback=teal!4,colframe=teal!35!black,title={Before you move on}]
What does \hi{झूला} mean? (\textquotedblleft{}A swing\textquotedblright{}.) Say it once more.
\end{tcolorbox}
\section[\texorpdfstring{\hi{पर्दा} (pardā)}{पर्दा (pardā)}]{\hi{पर्दा} (pardā) — a curtain}

\label{lesson:HI-C265-parda}



\begin{quote}
Before the new one: say the Hindi for a kite, then the Hindi for a swing.
\end{quote}

\subsection*{You'll want to know: \hi{पर्दा}}
\textbf{\hi{पर्दा}} — \emph{pardā} — \textquotedblleft{}a curtain\textquotedblright{}.

\begin{cousinweb}
\textbf{\hi{पर्दा}} (\emph{pardā}) is a loan, usually traced to Persian \emph{parda}, \textquotedblleft{}curtain, veil\textquotedblright{}.
\end{cousinweb}

\begin{grammarlens}[title={What you've built}]
One more word for this chapter.
\end{grammarlens}

\subsection*{Your turn}
\begin{itemize}
\raggedright
  \item \emph{Say it:} \emph{pardā}
  \item \emph{Say it:} \emph{pardā}, once more
  \item \emph{Say it:} say \emph{pardā}
  \item \emph{Recall:} say the Hindi for a kite, then the Hindi for a swing, then say \emph{pardā} again
\end{itemize}

\begin{tcolorbox}[breakable,colback=teal!4,colframe=teal!35!black,title={Before you move on}]
What does \hi{पर्दा} mean? (\textquotedblleft{}A curtain\textquotedblright{}.) Say it once more.
\end{tcolorbox}
\section[\texorpdfstring{\hi{दरी} (darī)}{दरी (darī)}]{\hi{दरी} (darī) — a cotton rug}

\label{lesson:HI-C265-dari}



\begin{quote}
Before the new one: say the Hindi for a swing, then the Hindi for a curtain.
\end{quote}

\subsection*{You'll want to know: \hi{दरी}}
\textbf{\hi{दरी}} — \emph{darī} — \textquotedblleft{}a cotton rug\textquotedblright{}.

\begin{cousinweb}
\textbf{\hi{दरी}} (\emph{darī}) is a tadbhava, a word inherited and reshaped by sound change. It comes from \emph{*darikā}, \textquotedblleft{}split cane, mat\textquotedblright{}, a form CDIAL reconstructs from its descendants (CDIAL 6194). CDIAL glosses the Hindi word \textquotedblleft{}carpet\textquotedblright{}.
\end{cousinweb}

\begin{grammarlens}[title={What you've built}]
One more word for this chapter.
\end{grammarlens}

\subsection*{Your turn}
\begin{itemize}
\raggedright
  \item \emph{Say it:} \emph{darī}
  \item \emph{Say it:} \emph{darī}, once more
  \item \emph{Say it:} say \emph{darī}
  \item \emph{Recall:} say the Hindi for a swing, then the Hindi for a curtain, then say \emph{darī} again
\end{itemize}

\begin{tcolorbox}[breakable,colback=teal!4,colframe=teal!35!black,title={Before you move on}]
What does \hi{दरी} mean? (\textquotedblleft{}A cotton rug\textquotedblright{}.) Say it once more.
\end{tcolorbox}
\section[\texorpdfstring{\hi{क़ालीन} (qālīn)}{क़ालीन (qālīn)}]{\hi{क़ालीन} (qālīn) — a carpet}

\label{lesson:HI-C265-qalin}



\begin{quote}
Before the new one: say the Hindi for a curtain, then the Hindi for a cotton rug.
\end{quote}

\subsection*{You'll want to know: \hi{क़ालीन}}
\textbf{\hi{क़ालीन}} — \emph{qālīn} — \textquotedblleft{}a carpet\textquotedblright{}.

\begin{cousinweb}
\textbf{\hi{क़ालीन}} (\emph{qālīn}) is a loan, usually traced to Persian \emph{qālī}, \textquotedblleft{}carpet\textquotedblright{}.
\end{cousinweb}

\begin{grammarlens}[title={What you've built}]
One more word for this chapter.
\end{grammarlens}

\subsection*{Your turn}
\begin{itemize}
\raggedright
  \item \emph{Say it:} \emph{qālīn}
  \item \emph{Say it:} \emph{qālīn}, once more
  \item \emph{Say it:} say \emph{qālīn}
  \item \emph{Recall:} say the Hindi for a curtain, then the Hindi for a cotton rug, then say \emph{qālīn} again
\end{itemize}

\begin{tcolorbox}[breakable,colback=teal!4,colframe=teal!35!black,title={Before you move on}]
What does \hi{क़ालीन} mean? (\textquotedblleft{}A carpet\textquotedblright{}.) Say it once more.
\end{tcolorbox}
